\scanpage{60}
% archive sheet number: 59
\uline{Lemme} $C$ une petite catégorie, $\pi$ \inserted{petit} cardinal, $E=\operatorname{Ind}_{\pi}(C)$ la sous-catégorie pleine de $\underline{\operatorname{Ind}}(C)$ formée des syst. inductifs \inserted{indexés par des ens. ordonnés} grands devant $\pi$. Alors

(i) $E$ \inserted{(est une $\mathcal U$-catégorie qui)} satisfait $L_\pi$, et l'inclusion $E\to\operatorname{Ind}(C)$ \deleted{commute} est $\pi$-accessible.

(ii) Les objets de $E$ qui sont $\pi$-accessibles sont ceux qui sont isom. à des facteurs directs d'objets de $C$ (donc forment une petite catégorie équivalente à l'enveloppe karoubienne de $C$, compte \deleted{tenu} que $\operatorname{Ind}_\pi(C)$ est stable par facteurs directs, cf \unsure{(i)$'$]|(i)]|(i))}. En particulier, si $C$ est karoubienne, \deleted{c'est que \unsure{l'appartenance d'un|les pts isolés} à $C$ alors $E$} \inserted{ce sont \unsure{préc.|pr.} ceux qui sont isom ds $E$ à un} objet de $C$.

(iii) $C$ est strictement génératrice dans $E$

(iv) Pour tt objet $X$ de $E$, $C/X$ est une \inserted{(petite)} catégorie filtrante grande devant $\pi$

\deleted{\unsure{NB Donc}}
\uline{Réciproquement :} Soit $E$ une \deleted{\unsure{$\mathcal U$-}}catégorie, \inserted{$\pi$ un \inserted{petit} cardinal}. Alors $E$ est $\simeq$ à une catégorie de la forme $\operatorname{Ind}_\pi(C)$, \inserted{$C$ petite,} \unsure{ssi il|ssi elle} satisfait les conditions suivantes

(i) $E$ \inserted{est une $\mathcal U$-catégorie et} satisfait $L_\pi$

(ii) La sous-catégorie \inserted{$E_\pi$} de $E$ formée des objets $\pi$-accessibles est \deleted{petit} petite, et strictement génératrice

(iii) Pour tt objet $X$ de $E$, la catégorie $E_{\pi/X}$ est filtrante et grande devant $\pi$.

\uline{NB} La condition (iii) est \deleted{vérifiée} \inserted{conséquence de (i),(ii)} si $E$ stable par \inserted{petites} $\varinjlim$ (alors $E_\pi$ est stable par $\varinjlim$ de cardinal $\leq\pi$)

\deleted{Soit $\pi'$ un cardinal, et}
Revenons à $\operatorname{Ind}_\pi(C)$, soit $\pi'$ un cardinal $>\pi$, et soit $\operatorname{Ind}_{\pi}(C)_{\pi'}$ \inserted{$=\operatorname{Fil}^{\pi'}(E)$} la sous-catégorie str. pleine de $\operatorname{Ind}_\pi(C)$ formée des ind-objets indexés par des $I$ tels que a) $I$ grand devant $\pi$ b) \deleted{\unsure{$\operatorname{card} I$|$\operatorname{card}\operatorname{Ob} I$}} \inserted{$\operatorname{Ob} I$} a un ensemble cofinal de cardinal $\leq\pi'$. Alors

a) Les objets de $\operatorname{Ind}_\pi(C)_{\pi'}$ \inserted{$\operatorname{Fil}^{\pi'}(E)$} dans $E$ sont \inserted{des objets} $\pi'$-accessibles :
\margin{en particulier, tt objet de $E$ est accessible !}% slanted
\margin{$\operatorname{Fil}^{\pi'}(E)\deleted{\unsure{=}}\subset E_{\pi'}$}

b) Si $\pi'$ est de la forme $2^c$ \inserted{($c\geq\pi$)}, alors la réciproque est vraie i.e. $\operatorname{Fil}^{\pi'}(E)=E_{\unsure{\pi'|\pi}}$
\margin{Peut-on avoir $\operatorname{Fil}^{\pi'}(E)\ne E_{\pi'}$ ?}% slanted

c) La filtration de $E$ par les \deleted{$\operatorname{Ind}_\pi(C)$ \unsure{ou par son}} $\operatorname{Fil}^{\pi'}(E)$ \inserted{($\pi'>\pi$)} est une ``filtration cardinale'' \ldots du moins si $E$ stable par petites $\varinjlim$, i.e. \inserted{$C$ stable par $\varinjlim$ de card $\leq\pi$}.
\margin{prouver c) avant b)}% slanted

\deleted{d') (Si $E$ est stable par petites $\varinjlim$) Pour tt cardinal $\pi'\geq\pi$, $E$ est équivalente à $\operatorname{Ind}(E_{\pi'})$. $(E_{\pi'})_{\pi'\geq\pi}$ forment une autre filtration cardinale.}
\deleted{$C$ stable par $\varinjlim$ de card \unsure{finie|$\leq\pi'$}}

\inserted{d) Si $E$ \deleted{\unsure{l'ind}} est stable par $\varinjlim$ \deleted{et} pour tt cardinal $\pi'\geq\pi$, $E\simeq\operatorname{Ind}(E_{\pi'})$.}

\margin{on a égalité si $E$ stable par $\varinjlim$}
\margin{cf SGA4 I 9.16 pour $\supset$}
\margin{résulte de la réciproque plus haut}
